
A manipulação de arquivos em C envolve operações como criação, abertura, 
leitura, escrita e fechamento de arquivos. Como discutido anteriormente na 
seção \ref{o-que-sao-arquivos}, o terminal também se enquadra na definição de 
arquivos, pois pode receber operações de escrita, leitura, entre outras.

Em aplicações mais complexas, pode ser necessário armazenar algumas informações 
fornecidas pelo usuário para reutilização pelo programa na próxima execução. 
Por exemplo, seria impraticável e inconveniente ter que criar uma nova conta 
ou reconfigurar todas as opções toda vez que um 
jogo é iniciado.

O uso da manipulação de arquivos é crucial devido às seguintes características:

\begin{itemize}
        \renewcommand{\labelitemi}{\tiny$\blacksquare$}
    \item \textbf{Reusabilidade:} As informações armazenadas em um arquivo 
        podem ser acessadas, atualizadas e removidas conforme necessário.
    \item \textbf{Portabilidade:} Como as informações são geralmente escritas 
        seguindo algum padrão lógico, um arquivo de configurações pode ser 
        facilmente transferido para outro computador sem complicações, 
        reduzindo erros durante a transferência de dados.
    \item \textbf{Eficiência:} Em muitas situações, é mais eficiente gerenciar 
        um grande volume de informações do usuário com apenas algumas 
        instruções, economizando tempo e reduzindo erros.
    \item \textbf{Capacidade de Armazenamento:} Armazenar informações em 
        arquivo elimina a necessidade de manter todas as informações de uma 
        aplicação na memória RAM simultaneamente, evitando preocupações e erros 
        relacionados à sobrecarga de memória.
\end{itemize}

\subsection{Tipos de arquivos}

\subsubsection{Arquivos de texto}

Esses arquivos contêm dados em formato de caracteres ASCII, podendo ser lidos e 
editados por qualquer editor de texto. Geralmente, mas não obrigatoriamente, 
esses arquivos possuem uma extensão '.txt' e todas as suas linhas terminam com 
o caractere '\textbackslash n'.

É comum encontrar esses tipos de arquivos sendo usados para guardar, por exemplo:

\begin{itemize}
        \renewcommand{\labelitemi}{\tiny$\blacksquare$}
    \item \textbf{Registro de logs:} Os arquivos de texto são comumente usados 
        para registrar eventos e informações relevantes durante a execução de 
        um programa. Por exemplo, um servidor web pode registrar todas as 
        solicitações de página em um arquivo de log para fins de monitoramento 
        e análise.
    \item \textbf{Listas de contatos:} Arquivos de texto também podem ser usados 
        para armazenar dados estruturados simples, como listas de contatos, 
        configurações de aplicativos, ou até mesmo informações de configuração 
        de rede.
    \item \textbf{Arquivos de código fonte:} Embora não seja exclusivo de 
        arquivos de texto, os arquivos de código-fonte em linguagens de 
        programação são geralmente armazenados em formato de texto para 
        facilidade de leitura e manipulação. Eles são editáveis por qualquer 
        editor de texto e podem ser versionados facilmente por sistemas de 
        controle de versão como Git.
\end{itemize}

\subsubsection{Arquivos binários}

Diferentemente dos arquivos de texto, que têm o seu conteúdo em formato de 
caracteres ASCII, os arquivos binários têm o seu conteúdo em um formato numérico, 
que fora de contexto não podem ser interpretados por um ser humano, mas que, 
quando corretamente interpretados por um programa, podem ser lidos e 
transformados em informação útil. Por conta disso, os arquivos binários tendem 
a ter uma maior segurança, pois mesmo que caiam em mãos de usuários maliciosos, 
não podem ser facilmente interpretados e manipulados de maneira mal intencionada.

Geralmente, esse tipo de arquivo é usado nas seguintes aplicações:

\begin{itemize}
        \renewcommand{\labelitemi}{\tiny$\blacksquare$}
    \item \textbf{Bancos de dados locais:} Muitos sistemas de gerenciamento de 
        banco de dados usam arquivos binários para armazenar dados localmente. 
        Embora os bancos de dados sejam mais comumente associados a servidores 
        e sistemas de gerenciamento de dados complexos, eles também podem ser 
        usados para aplicações simples de desktop.
    \item \textbf{Guardar imagens e multimídias:} Arquivos binários são 
        amplamente utilizados para armazenar imagens, vídeos, áudio e outros 
        tipos de multimídia. Isso ocorre porque esses tipos de dados são melhor 
        representados em formato binário, o que permite uma manipulação 
        eficiente e compactação.
    \item \textbf{Arquivos executáveis:} Os arquivos binários também incluem 
        executáveis de programas. Quando você compila um programa em linguagem 
        de máquina, o resultado é um arquivo binário que contém as instruções 
        que o computador pode executar diretamente.
\end{itemize}

\subsection{fopen, fread, fwrite e fclose}

\subsubsection{fopen}

A função \texttt{fopen} é usada para abrir um arquivo existente ou criar um 
novo arquivo. Ela retorna um ponteiro para o arquivo que pode ser usado em 
outras operações de entrada e saída de arquivo. A assinatura da função é a 
seguinte:

\begin{excode}
    FILE *fopen(const char *filename, const char *mode);
\end{excode}

O primeiro argumento \texttt{filename} é um vetor que contém o caminho do 
arquivo a ser aberto ou criado. O segundo argumento \texttt{mode} especifica o 
modo de abertura do arquivo, como leitura ("r"), escrita ("w"), anexação ("a"), 
etc.

\subsubsection{fread}

A função \texttt{fread} é usada para ler blocos de dados de um arquivo. Ela 
recebe quatro argumentos: um ponteiro para o vetor onde os dados lidos serão 
armazenados, o tamanho em bytes de cada elemento a ser lido, o número de 
elementos a serem lidos e um ponteiro para o arquivo a partir do qual os dados 
serão lidos. A assinatura da função é a seguinte:

\begin{excode}
    size_t fread(void *ptr, size_t size, size_t nmemb, FILE *stream);
\end{excode}

O retorno da função é o número total de elementos lidos com sucesso. Valor esse 
que pode ser utilizado para identificar quando um arquivo foi lido por completo, 
já que quando a função alcança o final do arquivo, 0 bytes são lidos do mesmo.

\subsubsection{fwrite}

A função \texttt{fwrite} é usada para escrever blocos de dados em um arquivo. 
Ela recebe quatro argumentos: um ponteiro para um vetor de dados a ser escrito, 
o tamanho em bytes de cada elemento a ser escrito, o número de elementos a 
serem escritos e um ponteiro para o arquivo no qual os dados serão escritos. A 
assinatura da função é a seguinte:

\begin{excode}
    size_t fwrite(const void *ptr, size_t size, size_t nmemb, FILE *stream);
\end{excode}

O retorno da função é o número total de elementos gravados com sucesso.

\subsubsection{fclose}

A função \texttt{fclose} é usada para fechar um arquivo previamente aberto. Ela 
recebe um ponteiro para o arquivo a ser fechado como seu único argumento. 
A assinatura da função é a seguinte:

\begin{excode}
    int fclose(FILE *stream);
\end{excode}

Se a operação for bem-sucedida, \texttt{fclose} retorna zero. Se ocorrer um 
erro, retorna um valor diferente de zero.

Essas funções são essenciais para manipular arquivos em C e são amplamente 
utilizadas em programas que necessitam de entrada e saída de dados persistentes.

\subsubsection{Exemplo prático}

Após entender o que cada uma das funções acima faz, vamos tentar exemplificar um uso prático das mesmas:

\begin{excode}
    #include <stdio.h>
    #include <stdlib.h>
    #include <ctype.h>

    int main() {
        // Abre o arquivo para leitura
        FILE *arquivo_entrada = fopen("dados.txt", "r");
        if (arquivo_entrada == NULL) {
            perror("Erro ao abrir arquivo de entrada");
            return EXIT_FAILURE;
        }

        // Abre um novo arquivo para escrita
        FILE *arquivo_saida = fopen("dados_modificados.txt", "w");
        if (arquivo_saida == NULL) {
            perror("Erro ao criar arquivo de saída");
            fclose(arquivo_entrada);
            return EXIT_FAILURE;
        }

        // Lê o conteúdo do arquivo de entrada
        char buffer[100];
        size_t bytes_lidos = fread(buffer, sizeof(char), 100, arquivo_entrada);
        if (bytes_lidos == 0) {
            perror("Erro ao ler arquivo de entrada");
            fclose(arquivo_entrada);
            fclose(arquivo_saida);
            return EXIT_FAILURE;
        }

        // Modifica os dados lidos (exemplo: converter para maiúsculas)
        for (size_t i = 0; i < bytes_lidos; i++) {
            buffer[i] = toupper(buffer[i]);
        }

        // Escreve os dados modificados no arquivo de saída
        size_t bytes_escritos = fwrite(buffer, sizeof(char), bytes_lidos, arquivo_saida);
        if (bytes_escritos != bytes_lidos) {
            perror("Erro ao escrever no arquivo de saída");
            fclose(arquivo_entrada);
            fclose(arquivo_saida);
            return EXIT_FAILURE;
        }

        // Fecha os arquivos abertos
        if (fclose(arquivo_entrada) == EOF) {
            perror("Erro ao fechar arquivo de entrada");
            return EXIT_FAILURE;
        }
        if (fclose(arquivo_saida) == EOF) {
            perror("Erro ao fechar arquivo de saída");
            return EXIT_FAILURE;
        }

        printf("Arquivo modificado com sucesso!\n");

        return EXIT_SUCCESS;
    }
\end{excode}

Neste exemplo:

\begin{itemize}
        \renewcommand{\labelitemi}{\tiny$\blacksquare$}
    \item Abrimos um arquivo chamado \texttt{dados.txt} para leitura ("r") e 
        outro chamado \\\texttt{dados\_modificados.txt} para escrita ("w").
    \item Usamos \texttt{fread} para ler até 100 bytes do arquivo de entrada.
    \item Modificamos os dados lidos (nesse caso, convertemos todas as letras para maiúsculas).
    \item Usamos \texttt{fwrite} para escrever os dados modificados no arquivo de saída.
    \item Finalmente, fechamos ambos os arquivos com \texttt{fclose}.
\end{itemize}

\begin{postscript}
    \begin{quote}
    "Dê um peixe a um homem e ele se alimentará por um dia. Ensine-o a pescar e ele se alimentará pelo resto da vida."
    \end{quote}

    Seria injusto que em um livro eu apenas ensinasse como usar as funções 
    \texttt{fopen}, \texttt{fread}, \texttt{fwrite} e \texttt{fclose}, sem 
    considerar a possibilidade de que novos padrões possam surgir no futuro, ou 
    que você precise atualizar seus conhecimentos. Por isso, vou compartilhar 
    um pequeno "truque" muito utilizado entre programadores em linguagem C, que 
    pode ser útil para relembrar ou aprender como uma função funciona.

    Para isso, vamos usar o comando Linux \texttt{man} (abreviação de manual), 
    cujo nome por si só é autoexplicativo. Ele é usado para pesquisar o manual 
    de utilização de alguma função ou comando. Felizmente, também é possível 
    pesquisar por funções da biblioteca padrão da linguagem C, conhecida como 
    \texttt{libc}. Para aqueles que não usam sistemas Linux, é possível usar o 
    comando \texttt{man} para pesquisar sobre funções até mesmo pelo Google. 

    Basta usar o comando:

    \begin{textbox}
        man nome_da_funcao_ou_comando
    \end{textbox}

    Com isso, você será capaz de encontrar o manual de qualquer função padrão 
    da biblioteca \texttt{libc}, e até mesmo de algumas funções que estão fora 
    das bibliotecas padrão.
\end{postscript}

\subsection{fseek, ftell e rewind}

Essa seção aborda funções específicas para manipulação do ponteiro de posição 
em arquivos, permitindo navegar por diferentes partes de um arquivo sem 
necessariamente ler todo o seu conteúdo sequencialmente.

\subsubsection{fseek}

A função \texttt{fseek()} possui o seguinte protótipo:

\begin{excode}
int fseek(FILE *stream, long offset, int whence);
\end{excode}

A função \texttt{fseek} move o ponteiro de leitura/escrita do arquivo para uma 
posição específica. Segue abaixo uma lista explicando os parâmetros da função:

\begin{itemize}
    \item \texttt{stream}: Ponteiro para o arquivo (obtido com \texttt{fopen()})
    \item \texttt{offset}: Número de bytes de deslocamento (pode ser positivo ou negativo)
    \item \texttt{whence}: Ponto de referência para o deslocamento que aceita os valores \texttt{SEEK\_SET}, \texttt{SEEK\_CUR} ou \texttt{SEEK\_END}.
\end{itemize}

Sendo \texttt{SEEK\_SET} o início do arquivo, \texttt{SEEK\_CUR} a posição atual do arquivo e \texttt{SEEK\_END} 
o final do arquivo. Além disso, a função retorna 0 em caso de sucesso e um valor diferente de 0 em caso de erro.

\subsubsection{Exemplos Práticos}

\begin{excode}
#include <stdio.h>

int main() {
    FILE *fp;
    
    // Abre o arquivo para leitura
    fp = fopen("dados.txt", "r");
    
    if (fp == NULL) {
        printf("Erro ao abrir o arquivo!\n");
        return 1;
    }
    
    // 1. Ir para o final do arquivo
    fseek(fp, 0, SEEK_END);
    
    // 2. Retornar ao início do arquivo
    fseek(fp, 0, SEEK_SET);
    
    // 3. Avançar 10 bytes a partir da posição atual
    fseek(fp, 10, SEEK_CUR);
    
    // 4. Retroceder 5 bytes a partir da posição atual
    fseek(fp, -5, SEEK_CUR);
    
    // 5. Posicionar 50 bytes antes do final do arquivo
    fseek(fp, -50, SEEK_END);
    
    fclose(fp);
    return 0;
}
\end{excode}

\subsubsection{ftell}

A função \texttt{ftell()} possui o seguinte protótipo:

\begin{excode}
long ftell(FILE *stream);
\end{excode}

A função funciona de forma bem simples, simplesmente retornando a posição atual do ponteiro no 
arquivo relativa ao seu início (em bytes). Sendo assim, extremamente conveniente para calcular o 
tamanho de um arquivo, como é mostrado no código abaixo:

\begin{excode}
#include <stdio.h>

int main() {
    FILE *fp = fopen("arquivo.txt", "r");
    long tamanho, pos_atual;
    
    if (fp == NULL) {
        printf("Erro ao abrir arquivo!\n");
        return 1;
    }
    
    // Vai para o final
    fseek(fp, 0, SEEK_END);
    
    // Obtém o tamanho do arquivo
    tamanho = ftell(fp);
    printf("Tamanho do arquivo: %ld bytes\n", tamanho);
    
    // Retorna ao início
    fseek(fp, 0, SEEK_SET);
    
    fclose(fp);
    return 0;
}
\end{excode}

\subsubsection{rewind}

A função \texttt{rewind()} possui o seguinte protótipo:

\begin{excode}
void rewind(FILE *stream);
\end{excode}

Essa função reposiciona o ponteiro para o início do arquivo, sendo assim equivalente a:

\begin{excode}
fseek(fp, 0, SEEK_SET);
\end{excode}

Com um único porém que o indicador de erro é limpo.  O que pode também ser feito da 
seguinte forma:

\begin{excode}
fseek(fp, 0, SEEK_SET);
clearerr(fp);
\end{excode}

Segue abaixo um exemplo de como utilizar a função \texttt{rewind}:

\begin{excode}
#include <stdio.h>

int main() {
    FILE *fp = fopen("exemplo.txt", "r+");
    char buffer[100];
    
    if (fp == NULL) {
        printf("Erro ao abrir arquivo!\n");
        return 1;
    }
    
    // Lê uma linha do arquivo
    fgets(buffer, 100, fp);
    printf("Primeira leitura: %s", buffer);
    
    // Retorna ao início para ler novamente
    rewind(fp);
    
    // Lê a mesma linha novamente
    fgets(buffer, 100, fp);
    printf("Leitura apos rewind: %s", buffer);
    
    fclose(fp);
    return 0;
}
\end{excode}

\begin{postscript}
    O indicador de erro de um arquivo é um \textbf{sinalizador interno} (flag) associado 
    ao ponteiro de arquivo \texttt{FILE} que indica se a última operação de entrada/saída 
    teve sucesso ou falhou. Quando uma função como \texttt{fread()}, \texttt{fwrite()}, 
    \texttt{fscanf()} falha, esse indicador é ativado.

    Portanto, ao limpar o indicador de erro, é possível tentar fazer novamente a parte do 
    processo que falhou tendo o benefício de poder usar o indicador de erro novamente.
\end{postscript}

\subsubsection{Considerações Importantes}

\begin{itemize}
    \item \textbf{Arquivos binários vs texto:} Em arquivos binários, \texttt{fseek()} com 
        \texttt{SEEK\_END} funciona normalmente. Em arquivos de texto, alguns sistemas 
        podem não suportar deslocamentos negativos.
    
    \item \textbf{Portabilidade:} Para arquivos muito grandes (mais de 2GB), use \texttt{fseeko()} 
        e \texttt{ftello()} que utilizam \texttt{off\_t} em vez de \texttt{long}.
    
    \item \textbf{Verificação de erros:} Sempre verifique o retorno de \texttt{fseek()} 
        em aplicações críticas:
    \begin{excode}
    if (fseek(fp, offset, SEEK_SET) != 0) {
        printf("Erro ao reposicionar ponteiro!\n");
    }
    \end{excode}
\end{itemize}
